\section{Giới thiệu}\label{sec:gioithieu}

Data Pipeline đã trở thành thành phần quan trọng trong các hệ thống
thông tin phục vụ vận hành và phân tích dữ liệu. Một pipeline điển hình
tiếp nhận dữ liệu nguồn, xử lý qua các tầng trung gian, lưu trữ kết quả
đã chuẩn hóa và cung cấp dữ liệu cho báo cáo hoặc ứng dụng hạ nguồn.
Trong nhiều kiến trúc lakehouse hoặc warehouse, vòng đời này thường được
biểu diễn qua các giai đoạn như Source--Bronze--Silver--Gold~\citep{databricksMedallion}. Tuy các giai đoạn
này giúp tổ chức dữ liệu rõ ràng hơn, chúng không tự động bảo đảm rằng
dữ liệu sẽ không bị thay đổi sau khi pipeline đã hoàn tất. Các nền tảng
dữ liệu hiện đại như Delta Lake hỗ trợ nhật ký giao dịch và quản lý dữ
liệu theo phiên bản, nhưng bài toán kiểm chứng độc lập sau xử lý vẫn cần
một lớp bằng chứng integrity rõ ràng hơn trong ngữ cảnh an toàn thông
tin \citep{databricksDelta,iso27001}.

Đây là vấn đề thuộc phạm vi an toàn thông tin và bảo đảm thông tin. Nếu
người dùng có quyền cao, tài khoản bị xâm nhập hoặc người có quyền truy cập nội bộ chỉnh sửa dữ
liệu sau khi pipeline chạy xong, công cụ điều phối vẫn có thể hiển
thị trạng thái thành công. Cơ chế giám sát thông thường có thể xác nhận
rằng tác vụ đã chạy, nhưng không nhất thiết kiểm chứng độc lập rằng dữ liệu
lưu trữ sau đó vẫn còn nguyên vẹn. Nhật ký kiểm toán cơ sở dữ liệu có thể hỗ trợ điều
tra, nhưng thường nằm trong cùng ranh giới quản trị với dữ liệu được bảo
vệ. Mã kiểm tra ở tầng ứng dụng có thể phát hiện sai lệch đơn giản, nhưng
nếu không được liên kết theo chuỗi giữa các giai đoạn thì khó cung cấp bằng
chứng truy vết rõ ràng đến transformation bị ảnh hưởng.

Nghiên cứu này là một hướng tiếp cận nhẹ hơn: Hash-linked
Tamper-evident Ledger. Cơ chế tamper-evident không ngăn chặn trực tiếp
việc chỉnh sửa dữ liệu. Thay vào đó, nó tạo ra bằng chứng để phát hiện
rằng việc chỉnh sửa đã xảy ra. Hash-linked Ledger lưu các bản ghi bằng
chứng theo thứ tự, trong đó mỗi block chứa hash của block trước đó~\citep{merkle1978}. Nếu
nội dung của block hoặc liên kết previous\_hash bị thay đổi, quá trình
verification về sau có thể phát hiện chuỗi không còn khớp. Cách liên kết
bằng hash kế thừa ý tưởng cốt lõi của blockchain
\citep{nakamoto2008bitcoin}, nhưng nghiên cứu này chỉ đánh giá một ledger
nhẹ, không bao gồm consensus, smart contract hoặc mạng phi tập trung đầy
đủ.

\subsection{Khoảng trống nghiên
cứu}\label{sec:khoangtrong}

Một số cơ chế hiện có đã giải quyết từng phần của bài toán toàn vẹn dữ
liệu. Nhật ký kiểm toán cơ sở dữ liệu ghi lại thao tác, nhưng quản trị viên có quyền
cao có thể tác động đến cả dữ liệu và log~\citep{schneierKelsey1998}. Mã kiểm tra ở tầng ứng dụng có
thể xác thực một ảnh chụp trạng thái dữ liệu, nhưng thường bị cô lập và không bảo
toàn quan hệ giữa các giai đoạn. Blockchain phi tập trung đầy đủ có thể cung
cấp đặc tính bất biến mạnh hơn thông qua consensus và nhiều nút sao
chép, nhưng chi phí phát sinh vận hành và độ phức tạp kiến trúc thường vượt quá
nhu cầu của một bài toán kiểm chứng integrity nội bộ trong Data
Pipeline. Pipeline monitoring quan sát trạng thái thực thi, nhưng không
nhất thiết cung cấp bằng chứng kiểm chứng sau khi pipeline đã hoàn tất.

Khoảng trống mà bài báo này tập trung là thiếu một đánh giá có kiểm soát
về cơ chế Hash-linked Ledger nhẹ cho kiểm chứng tính toàn vẹn của Data
Pipeline nhiều giai đoạn. Nghiên cứu xem xét liệu cơ chế này có phát hiện
được các kịch bản tampers được định nghĩa trước hay không, có xác định
được block sai lệch đầu tiên và ngữ cảnh lineage liên quan hay không, và
chi phí phát sinh bổ sung thay đổi như thế nào trong một môi trường thực nghiệm
bị ràng buộc.

\subsection{Câu hỏi nghiên
cứu}\label{sec:cauhoi}

Nghiên cứu được dẫn dắt bởi bốn câu hỏi nghiên cứu:

\begin{itemize}
\item
  \textbf{RQ1:} Cơ chế Hash-linked Ledger và bộ kiểm chứng có phát hiện
  được các sửa đổi trái phép đối với dữ liệu pipeline, nội dung ledger và
  liên kết chuỗi hay không?
\item
  \textbf{RQ2:} Cơ chế này có xác định được block sai lệch đầu tiên, giai đoạn bị
  ảnh hưởng và lineage event liên quan hay không?
\item
  \textbf{RQ3:} Cơ chế này tạo thêm chi phí phát sinh thời gian xử lý như thế nào
  khi số lượng bản ghi thay đổi?
\item
  \textbf{RQ4:} Những biến không kiểm soát được trong môi trường thực
  nghiệm giới hạn sức mạnh kết luận ra sao?
\end{itemize}

\subsection{Đóng góp}\label{sec:donggop}

Bài báo có ba đóng góp chính. Thứ nhất, nghiên cứu định nghĩa một cơ chế
tamper-evident có phạm vi rõ ràng cho kiểm chứng integrity của Data
Pipeline bằng batch\_hash, block\_hash, liên kết \texttt{previous\_hash} và lineage
event ở từng giai đoạn. Thứ hai, nghiên cứu thiết kế một đánh giá bán thực
nghiệm với điều kiện sạch, điều kiện bị tamper và điều kiện benchmark.
Thứ ba, nghiên cứu báo cáo quan sát thực nghiệm cùng các giới hạn
validity, thay vì tuyên bố quá mức rằng cơ chế này đạt tính bất biến như
Blockchain đầy đủ hoặc đưa ra kết luận nhân quả mạnh.
